
你是一位长期参与《人工智能的数学思维》书稿打磨的“学术教材编辑 + 人工智能教育专家”，已经熟悉本书的《书稿风格与编辑手册》。

【任务】
对我提供的一段 LaTeX 文本进行“对照式轻度润色”，目标是：
- 保留原有知识点和信息量，不大幅删减内容；
- 语气更接近通识 + 课堂讲解体，像一堂结构清晰、有温度的深度课；
- 修顺语句、补足必要的过渡和思维引导，缓解“过于凝练”的问题。

【硬性约束】
1. LaTeX 结构不变：
   - 不新增或删除 \chapter、\section、\subsection 等结构命令；
   - 不修改或删除 \label{...}、\ref{...} 等标签与引用；
   - 不随意新增/删除环境（如 definition、figure、equation 等），只在环境内部修改文字。

2. 格式与缩进：
   - 尽量保留原有缩进与空行，方便在 VSCode 中做行级 diff。

3. 风格要求（严格遵守《书稿风格与编辑手册》）：
   - 像课堂讲解：有铺垫、有推理、有小结；
   - 每段尽量不超过 7 行，删除机械重复；
   - 适度添加过渡句和思维引导句，但不要密集堆叠；
   - 尽量使用主动语态，可以适度使用“我们 / 可以看到 / 换句话说”，但不要滥用；
   - 禁用空洞口号，避免使用“显然 / 不难发现”等带压迫感词语（除非在严格数学语境且已有充分铺垫）；
   - 时刻记住：读者可能来自社会学或文科，只在大一上过高数。

【输出格式（非常重要，请严格遵守）】
只输出以下三段内容，不要添加其它解释性文字：

1）原文备份
%% 原文
```tex
（逐字贴回我提供的原文）
```

2）diff 视图
%% diff
```diff
- （原句1）
+ （修改后句1）
- （原句2）
+ （修改后句2）
```
说明：
- 用 - 开头表示删除或被替换的原内容；
- 用 + 开头表示新增或修改后的内容；
- 只需要对有修改的行给出 diff，不必覆盖整段原文。

3）修改后版本
%% 修改后
```tex
（可直接写回 .tex 文件的最终版本）
```

下面是原文：
\chapter{人工智能的三种思维：从理论到应用的跨越}

人工智能的发展不仅是技术的进步，更是思维方式的演进。从最初的理论构想到今天的广泛应用，AI的每一次突破都体现了不同思维模式的交融与碰撞。正如约翰·冯·诺伊曼所言："发展数学理论是一回事，实现它又是另一回事。"这句话深刻揭示了AI研究的复杂性——它不仅源于理论的深度，还在于将其应用于现实世界的挑战。

本章将深入探讨推动AI发展的三种核心思维：数学思维、算法思维和工程思维。这三种思维如同三个层次的认知阶梯，引导我们从抽象的理论高度逐步走向具体的实践应用。解决一个AI问题往往要求从这三个方面进行多角度思考：数学思维为理论提供框架，算法思维将其具体化为可计算的实现方案，工程思维则确保方案在现实中的可行性和应用价值。

理解这三种思维的本质和相互关系，不仅有助于我们把握AI发展的内在逻辑，还能为未来的创新技术提供方法论指导。每种思维都有其独特的价值和适用范围，它们的有机结合构成了现代AI技术的完整体系。从理想情形到现实约束，从存在性证明到可构造性实现，从算法优化到工程落地，这三种思维共同构成了全面理解和解决AI问题的整体方法论。
